\documentclass[11pt]{article}

\usepackage[margin=1in]{geometry}
\usepackage{amsmath, amssymb}
\usepackage{graphicx}
\usepackage{hyperref}
\usepackage{booktabs}

\title{Technical Assessment: Davis AI - Discrete Diffusion for Floorplans}
\author{Ahmad Fraij}
\date{\today}

\begin{document}

% ---------------- TITLE PAGE ----------------
\begin{titlepage}
    \centering
    \vspace*{2cm}

    {\LARGE \textbf{Technical Assessment}}\\[0.5cm]
    {\Large Davis AI - Discrete Diffusion for Floorplans}\\[2cm]

    {\large \textbf{Author:} Ahmad Fraij}\\[0.5cm]
    {\large \textbf{Date:} \today}\\[3cm]

    \vfill

\end{titlepage}

% ---------------- MAIN CONTENT ----------------

\section{Introduction}

The Graph2Plan pipeline generates floorplans from a building boundary and user constraints. It first retrieves layout graphs from the RPLAN dataset, then passes them into the Graph2Plan GNN to produce a rendered floor plan. The retrieval stage has problems with output quality, diversity, and memory usage that discrete diffusion can address.

\section{Issues with the Retrieve and Adjust Stage}

The retrieve and adjust stage does not generate layout graphs; it finds the closest boundary shape in the training set, extracts that sample's graph, then scales and shifts it to fit the new boundary. If no training example has a similar shape, the retrieved graph is a poor structural match with wrong room types, adjacencies, or room counts. These errors propagate through the GNN and degrade the final floor plan quality.

The other problem is memory. The retrieval stage keeps all training floor plan objects in memory at runtime, including boxes, edges, and boundaries for roughly 75K samples. This adds several GB of overhead that scales with dataset size.

\section{Graph Representation}

\subsection{Node Representation}

Each node represents one room and is encoded as a 48-dimensional vector made up of three concatenated one-hot parts:

\begin{itemize}
    \item \textbf{Room type} (13 dims): one of 13 room categories — LivingRoom, MasterRoom, Kitchen, Bathroom, DiningRoom, ChildRoom, StudyRoom, SecondRoom, GuestRoom, Balcony, Entrance, Storage, Wall.
    \item \textbf{Location} (25 dims): the room centre mapped onto a $5 \times 5$ grid over the boundary bounding box, encoded as a one-hot cell index.
    \item \textbf{Size} (10 dims): the room area as a fraction of the total boundary area, discretised into 10 equal-width bins.
\end{itemize}

\subsection{Edge Representation}

Edges are undirected. Each edge is a 2-dimensional one-hot vector: class 0 means no adjacency and class 1 means the two rooms are spatially adjacent. The graph is represented in dense format ($n \times n$ adjacency) during diffusion and converted to a sparse edge list for the GNN.

\section{Discrete Diffusion Formulation}

The model is based on DiGress~\cite{vignac2022digress}, a discrete denoising diffusion model for graphs.

\textbf{Forward process.} At each timestep $t \in \{1, \ldots, T\}$ noise is added to the clean graph $(X_0, E_0)$ by applying marginal transition matrices:
\[
q(\mathbf{z}_t \mid \mathbf{x}) = \mathbf{x} \, \bar{Q}_t
\]
where $\bar{Q}_t$ mixes the clean one-hot state with the data marginal distribution. The stationary distribution at $t = T$ is the empirical marginal over room types and edge types in the training set. A cosine noise schedule controls how fast the clean signal is corrupted across $T = 500$ steps.

\textbf{Reverse process.} A Graph Transformer network with 6 layers predicts the clean graph $\hat{x}_0$ from the noisy graph $\mathbf{z}_t$ and the conditioning vector $\mathbf{y}$. The posterior
\[
q(\mathbf{z}_{s} \mid \mathbf{z}_t, \mathbf{x}_0) \propto q(\mathbf{z}_t \mid \mathbf{z}_s)\, q(\mathbf{z}_s \mid \mathbf{x}_0)
\]
is then used to sample $\mathbf{z}_{t-1}$ from the prediction, stepping back to $t=0$.

\textbf{Training objective.} The simplified training loss is cross-entropy between the predicted clean graph and the ground truth:
\[
\mathcal{L} = \mathcal{L}_X + \lambda_E \mathcal{L}_E
\]
with $\lambda_E = 5$ upweighting the edge loss since edges are sparser than nodes. The label for node cross-entropy is taken as the argmax of the 48-dim target, which selects the room type.

\section{Conditioning Setup}

The model is conditioned on a 1164-dimensional vector $\mathbf{y}$ appended to every graph-level feature:

\[
\mathbf{y} = [\underbrace{\text{TF}}_{1000} \mid \underbrace{\text{counts}}_{14} \mid \underbrace{\text{location masks}}_{125} \mid \underbrace{\text{adjacency}}_{25}]
\]

\begin{itemize}
    \item \textbf{Turning Function (TF)} descriptor: the boundary polygon is encoded as a 1000-point piecewise-constant turning function that captures the shape regardless of rotation or scale.
    \item \textbf{Room counts} (14 dims): the required number of each of the 13 room types plus one bedroom-cluster count.
    \item \textbf{Location masks} (125 dims): for each of 5 constrained room types (Living, Master, Kitchen, Bathroom, Dining), a 25-bit mask marks which $5 \times 5$ grid cells the room may occupy.
    \item \textbf{Adjacency constraints} (25 dims): a $5 \times 5$ binary matrix indicating which pairs of constrained room types must be adjacent.
\end{itemize}

\textbf{Classifier-free guidance (CFG).} During training, 15\% of batches have the constraint portion of $\mathbf{y}$ (indices 1000 onward) zeroed out while keeping the TF descriptor intact. This lets the model learn both constrained and unconstrained generation, so at inference time constraints can be included or omitted.

\section{Benchmark Protocol}

\textbf{Dataset.} RPLAN~\cite{wu2019data}, a dataset of approximately 75,000 annotated residential floor plans. The split follows Graph2Plan: 70\% train, 15\% validation, 15\% test, using a fixed random seed for reproducibility.

\textbf{Baseline.} The original Graph2Plan retrieve-and-adjust stage, which copies the layout graph from the nearest training sample by boundary shape.

\textbf{Evaluation.} For each test sample, the baseline and the diffusion model each produce a room-type adjacency graph. Both graphs are fed into the same frozen Graph2Plan GNN to render a 128$\times$128 floor plan image. Metrics are computed against the ground-truth top-1 training reference for that test sample.

\textbf{Metrics.}
\begin{itemize}
    \item \textbf{FID / KID} ($\downarrow$): image-level distribution distance on rendered floor plans.
    \item \textbf{Connectivity} ($\uparrow$): fraction of generated graphs that are fully connected.
    \item \textbf{Graph validity} ($\uparrow$): fraction of graphs that are connected, contain at least one LivingRoom and one Bathroom, have at least one edge, and have no isolated nodes.
    \item \textbf{Room-type accuracy} ($\uparrow$): intersection-over-GT of predicted vs.\ ground-truth room-type multiset.
    \item \textbf{Room-type F1} ($\uparrow$): harmonic mean of room-type precision and recall.
    \item \textbf{Node count overlap} ($\uparrow$): $\min(n_\text{pred}, n_\text{gt}) / \max(n_\text{pred}, n_\text{gt})$.
    \item \textbf{Count satisfaction} ($\uparrow$): fraction of constrained room types whose generated count matches the target exactly.
    \item \textbf{Adjacency satisfaction} ($\uparrow$): fraction of required adjacency pairs present in the output.
    \item \textbf{Pixel mIoU} ($\uparrow$): mean per-class pixel IoU between the rendered output and ground-truth segmentation.
\end{itemize}

\section{Results}

\begin{table}[h]
\centering
\begin{tabular}{lcccccc}
\toprule
Model & FID $\downarrow$ & KID $\downarrow$ & Validity $\uparrow$ & RT-F1 $\uparrow$ & Count Sat. $\uparrow$ & mIoU $\uparrow$ \\
\midrule
Retrieval baseline & -- & -- & -- & -- & -- & -- \\
Ours (constrained diffusion) & -- & -- & -- & -- & -- & -- \\
\bottomrule
\end{tabular}
\caption{Comparison on RPLAN test split (500 samples). Results to be filled after full evaluation run.}
\end{table}

\section{Ablation Studies}

Two design choices were investigated:

\textbf{Constraint dropout rate.} The CFG dropout probability was set to 0.15. A higher rate makes the unconstrained mode stronger but weakens constraint following; a lower rate hurts unconditional generation quality. 0.15 was chosen as a reasonable default following standard CFG practice.

\textbf{Edge loss weight $\lambda_E$.} With $\lambda_E = 1$ the model tends to generate sparse graphs with missing adjacencies. Setting $\lambda_E = 5$ improves edge recall without visibly hurting node type accuracy, since there are far fewer edge classes (2) than node classes (48).

\section{Failure Cases}

The main failure mode is graphs where a room is disconnected from the rest. This happens when the model places many rooms but generates too few edges, particularly for large room counts sampled from the node count distribution. The post-processing step catches this by connecting isolated components through the LivingRoom, but the forced edges are structurally arbitrary.

A second issue is room type imbalance. Rare types like StudyRoom and SecondRoom are underrepresented in the training data, so the model under-generates them even when they are specified in the constraint vector.

\section{Future Work}

The multi-hot node feature representation (room type + location + size in 48 dims) does not form a valid one-hot categorical distribution, which technically breaks the discrete diffusion assumptions for the location and size components. A cleaner formulation would either run three separate diffusion channels per node or use a joint class index over the combined discrete space. This is the most important technical fix to address before scaling up training.

On the conditioning side, location masks and adjacency constraints are currently set to zero at test time since they are not known beforehand. Using predicted or user-supplied location constraints rather than zeros would likely improve count and adjacency satisfaction scores.

\section{Conclusion}

Replacing the retrieve-and-adjust stage with a discrete diffusion model addresses the core quality and diversity problems of the original pipeline. The model is conditioned on boundary shape and optional room constraints via classifier-free guidance, allowing flexible generation with or without explicit constraints. The main remaining technical issue is the multi-hot node feature representation, which should be resolved before extended training runs.

\begin{thebibliography}{9}

\bibitem{vignac2022digress}
C.\ Vignac, I.\ Krawczuk, A.\ Siraudin, B.\ Wang, V.\ Cevher, P.\ Frossard,
\textit{DiGress: Discrete Denoising Diffusion for Graph Generation},
arXiv:2209.14734, 2022.

\bibitem{hu2020graph2plan}
R.\ Hu, Z.\ Huang, Y.\ Tang, O.\ Van Kaick, H.\ Zhang, H.\ Huang,
\textit{Graph2Plan: Learning Floorplan Generation from Layout Graphs},
ACM Transactions on Graphics, 2020.

\bibitem{wu2019data}
W.\ Wu, X.\ Fu, R.\ Tang, Y.\ Wang, Y.\ Qi, W.\ Liu,
\textit{Data-driven Interior Plan Generation for Residential Buildings},
ACM Transactions on Graphics, 2019.

\end{thebibliography}

\end{document}
